% GENERATED FILE. Edit canonical lessons, then run npm run generate:books.
% canonical-source-hash: fnv1a64:19ba1a60048f539a
% canonical-lessons: PT-C486-termometro, PT-C486-vacina, PT-C486-injecao, PT-C486-entorse, PT-C486-ligadura

\chapter{Vacinas e ligaduras}
\label{ch:pt-vacinas-e-ligaduras}
\hlchaptermodality{486}

\begin{chapteropening}
\textbf{By the end of this chapter:} \emph{I can talk about thermometers, jabs, sprains and bandages.}

Complete the last lesson of chapter 486: I can talk about thermometers, jabs, sprains and bandages.
\end{chapteropening}

\section[o termómetro]{o termómetro — thermometer}

\label{lesson:PT-C486-termometro}



\begin{quote}
Before the new one: say the Portuguese for to swell, then the Portuguese for swollen.
\end{quote}

\subsection*{You'll want to know: o termómetro}
\textbf{o termómetro} — \textquotedblleft{}thermometer\textquotedblright{}.

\begin{quote}
O termómetro marcava trinta e nove graus.
\end{quote}

\begin{quote}
\emph{The thermometer read thirty-nine degrees.}
\end{quote}

In Brazil: \emph{termômetro}.

\begin{cousinweb}
\emph{termómetro} is usually traced to French \emph{thermomètre}, a learned coinage from Greek \emph{thermós} \textquotedblleft{}hot\textquotedblright{} and \emph{métron} \textquotedblleft{}measure\textquotedblright{}.
\end{cousinweb}

\begin{grammarlens}[title={What you've built}]
One more word for telling, explaining and describing at B1.
\end{grammarlens}

\subsection*{Your turn}
\begin{itemize}
\raggedright
  \item \emph{Say it:} \emph{o termómetro}
  \item \emph{Say it:} \emph{o termómetro}, once more
  \item \emph{Say it:} a sentence of your own with \emph{o termómetro}
  \item \emph{Recall:} say the Portuguese for to swell, then the Portuguese for swollen, then say \emph{o termómetro} again
\end{itemize}

\begin{tcolorbox}[breakable,colback=teal!4,colframe=teal!35!black,title={Before you move on}]
What does o termómetro mean? (\textquotedblleft{}Thermometer\textquotedblright{}.) Say it once more.
\end{tcolorbox}
\section[a vacina]{a vacina — vaccine, jab}

\label{lesson:PT-C486-vacina}



\begin{quote}
Before the new one: say the Portuguese for swollen, then the Portuguese for thermometer.
\end{quote}

\subsection*{You'll want to know: a vacina}
\textbf{a vacina} — \textquotedblleft{}vaccine, jab\textquotedblright{}.

\begin{quote}
As crianças levaram a vacina da gripe.
\end{quote}

\begin{quote}
\emph{The children had the flu jab.}
\end{quote}

\emph{Levar a vacina} — to get vaccinated.

\begin{cousinweb}
\emph{vacina} is a learned borrowing from Latin \emph{vaccīna} (EtymDB), from \emph{vaccīnus} \textquotedblleft{}of cows\textquotedblright{}, built on \emph{vacca} \textquotedblleft{}cow\textquotedblright{}: the first vaccine came from cowpox.
\end{cousinweb}

\begin{grammarlens}[title={What you've built}]
One more word for telling, explaining and describing at B1.
\end{grammarlens}

\subsection*{Your turn}
\begin{itemize}
\raggedright
  \item \emph{Say it:} \emph{a vacina}
  \item \emph{Say it:} \emph{a vacina}, once more
  \item \emph{Say it:} a sentence of your own with \emph{a vacina}
  \item \emph{Recall:} say the Portuguese for swollen, then the Portuguese for thermometer, then say \emph{a vacina} again
\end{itemize}

\begin{tcolorbox}[breakable,colback=teal!4,colframe=teal!35!black,title={Before you move on}]
What does a vacina mean? (\textquotedblleft{}Vaccine, jab\textquotedblright{}.) Say it once more.
\end{tcolorbox}
\section[a injeção]{a injeção — injection}

\label{lesson:PT-C486-injecao}



\begin{quote}
Before the new one: say the Portuguese for thermometer, then the Portuguese for vaccine.
\end{quote}

\subsection*{You'll want to know: a injeção}
\textbf{a injeção} — \textquotedblleft{}injection\textquotedblright{}.

\begin{quote}
A enfermeira deu-me uma injeção no braço.
\end{quote}

\begin{quote}
\emph{The nurse gave me an injection in the arm.}
\end{quote}

Before AO90 spelled \emph{injecção} in Portugal.

\begin{cousinweb}
\emph{injeção} is built on \emph{injetar} \textquotedblleft{}to inject\textquotedblright{}, a learned borrowing from Latin \emph{iniectāre} (EtymDB), from \emph{iacere} \textquotedblleft{}to throw\textquotedblright{}.
\end{cousinweb}

\begin{grammarlens}[title={What you've built}]
One more word for telling, explaining and describing at B1.
\end{grammarlens}

\subsection*{Your turn}
\begin{itemize}
\raggedright
  \item \emph{Say it:} \emph{a injeção}
  \item \emph{Say it:} \emph{a injeção}, once more
  \item \emph{Say it:} a sentence of your own with \emph{a injeção}
  \item \emph{Recall:} say the Portuguese for thermometer, then the Portuguese for vaccine, then say \emph{a injeção} again
\end{itemize}

\begin{tcolorbox}[breakable,colback=teal!4,colframe=teal!35!black,title={Before you move on}]
What does a injeção mean? (\textquotedblleft{}Injection\textquotedblright{}.) Say it once more.
\end{tcolorbox}
\section[a entorse]{a entorse — sprain}

\label{lesson:PT-C486-entorse}



\begin{quote}
Before the new one: say the Portuguese for vaccine, then the Portuguese for injection.
\end{quote}

\subsection*{You'll want to know: a entorse}
\textbf{a entorse} — \textquotedblleft{}sprain\textquotedblright{}.

\begin{quote}
Fiz uma entorse no tornozelo a jogar à bola.
\end{quote}

\begin{quote}
\emph{I sprained my ankle playing football.}
\end{quote}

Feminine in Portugal: \emph{uma entorse}.

\begin{cousinweb}
\emph{entorse} is usually traced to French \emph{entorse} \textquotedblleft{}a twist\textquotedblright{}, built on a verb that goes back to Latin \emph{torquēre} \textquotedblleft{}to twist\textquotedblright{} (Lewis \& Short).
\end{cousinweb}

\begin{grammarlens}[title={What you've built}]
One more word for telling, explaining and describing at B1.
\end{grammarlens}

\subsection*{Your turn}
\begin{itemize}
\raggedright
  \item \emph{Say it:} \emph{a entorse}
  \item \emph{Say it:} \emph{a entorse}, once more
  \item \emph{Say it:} a sentence of your own with \emph{a entorse}
  \item \emph{Recall:} say the Portuguese for vaccine, then the Portuguese for injection, then say \emph{a entorse} again
\end{itemize}

\begin{tcolorbox}[breakable,colback=teal!4,colframe=teal!35!black,title={Before you move on}]
What does a entorse mean? (\textquotedblleft{}Sprain\textquotedblright{}.) Say it once more.
\end{tcolorbox}
\section[a ligadura]{a ligadura — bandage}

\label{lesson:PT-C486-ligadura}



\begin{quote}
Before the new one: say the Portuguese for injection, then the Portuguese for sprain.
\end{quote}

\subsection*{You'll want to know: a ligadura}
\textbf{a ligadura} — \textquotedblleft{}bandage\textquotedblright{}.

\begin{quote}
O médico pôs-me uma ligadura no pulso.
\end{quote}

\begin{quote}
\emph{The doctor put a bandage on my wrist.}
\end{quote}

In Brazil: \emph{a atadura}.

\begin{cousinweb}
\emph{ligadura} is built on \emph{ligar} \textquotedblleft{}to tie\textquotedblright{}, which is inherited from Latin \emph{ligāre} \textquotedblleft{}to bind\textquotedblright{} (Lewis \& Short).
\end{cousinweb}

\begin{grammarlens}[title={What you've built}]
One more word for telling, explaining and describing at B1.
\end{grammarlens}

\subsection*{Your turn}
\begin{itemize}
\raggedright
  \item \emph{Say it:} \emph{a ligadura}
  \item \emph{Say it:} \emph{a ligadura}, once more
  \item \emph{Say it:} a sentence of your own with \emph{a ligadura}
  \item \emph{Recall:} say the Portuguese for injection, then the Portuguese for sprain, then say \emph{a ligadura} again
\end{itemize}

\begin{tcolorbox}[breakable,colback=teal!4,colframe=teal!35!black,title={Before you move on}]
What does a ligadura mean? (\textquotedblleft{}Bandage\textquotedblright{}.) Say it once more.
\end{tcolorbox}
